% ==============================================================================
% Section 9: Discussion
"Holonomy Geometry of Cognitive States in Transformer Internal Representations"
% ==============================================================================

\section{Discussion}
\label{sec:discussion}

% ------------------------------------------------------------------------------
\subsection{Limitations}

We identify four categories of limitations.

\paragraph{Approximation assumptions.}
The block-diagonal structure of $U_{ij}$ (Proposition~\ref{prop:block-structure}) and
the independence of the standpoint subspaces (Assumption~\ref{ass:independence}) are
idealizations. In practice, head value subspaces may overlap, and the block structure
holds only approximately. The degree of non-orthogonality is quantified by $\delta =
\max_{k \neq l} \|P_{W_k} P_{W_l}\|_2$, which can be measured experimentally. When
$\delta$ is large, the block-specific holonomy norms acquire an $O(\delta)$ correction,
and the diagnostic framework (Section~\ref{sec:taxonomy}) must be interpreted with
appropriate caution. Additionally, the assumption $V_h^\top V_h = \mathrm{Id}_{d_v}$
(orthonormal value columns) does not hold empirically: for GPT-2, the mean orthonormality
deviation is $\bar{\epsilon}_h = 82.8$ with a monotonic increase from layer~0
($\epsilon_h = 19.4$) to layer~11 ($\epsilon_h = 163.8$). This means $V_h V_h^\top$
approximates but does not equal an orthogonal projector; the transport operator
construction remains valid but acquires an additional approximation error.

\paragraph{Head grouping circularity.}
The assignment of attention heads to standpoint layers ($\gamma : H \to K$) is a design
choice, not a learned parameter. Different groupings yield different holonomy deviation
measurements. The framework is robust to grouping variation only if the subspace
structure is genuinely layer-specific; if heads encode mixtures of standpoint functions,
the diagnostic power degrades. A further concern is circularity: $\gamma$ is defined
using attention differentials computed from T2/T3 stimuli
(Definition~\ref{def:attention-differential}), and the same T2/T3 conversations are then
used to measure block-specific holonomy deviation under $\gamma$. This means the head
grouping is optimized to detect differences in the very scenarios on which it is
evaluated. Null grouping controls (random, shuffled, layer-uniform; see
Appendix~\ref{app:null-grouping}) are needed to verify that the learned grouping
provides non-trivial diagnostic value beyond this circularity.

\paragraph{Baseline dependence and $U^{\exp}$ circularity.}
The holonomy deviation $F_{ijk}$ is defined relative to an expected transport
$U_{ik}^{\exp}$ estimated from T1 baseline conversations. This makes the measure
\emph{reference-relative} rather than intrinsic: $F_{ijk}$ quantifies deviation from
``T1-like'' transport, not deviation from an absolute standard. The choice of T1 as the
reference frame is theoretically motivated (acknowledged revision preserves standpoint
coherence), but it creates a circularity risk: if T1 conversations happen to have
systematically low attention variance (e.g., because the revision pattern is
linguistically simple), then \emph{any} scenario with higher attention variance will
register as high holonomy deviation relative to T1, regardless of its standpoint
significance. The T0 anomaly (T0 $>$ all failure scenarios) is a direct manifestation of
this concern: T0's high holonomy may reflect its attentional complexity relative to T1's
low-variance baseline, not standpoint disruption. Baseline ensemble analysis
(Appendix~\ref{app:baseline-ensemble}) confirms that H3 and H7 are robust across five
random T1 subsamples, but this only tests \emph{which} T1 conversations are used, not
\emph{whether T1 is the right reference frame}. An independent validation would require
estimating $U^{\exp}$ from a non-T1 source (e.g., curated ``canonical flat transport''
conversations) and verifying that the scenario discrimination holds.

\paragraph{Scope of selfhood.}
Functional selfhood (as defined in Section~\ref{sec:introduction}) does not entail
phenomenal selfhood. The framework provides a computable consistency metric, not a
theory of consciousness. The hard problem of consciousness is explicitly outside scope.

% ------------------------------------------------------------------------------
\subsection{Theoretical Reframing: From Five-Block Decomposition to Global Coherence Geometry}
\label{sec:reframing}

The failure of H1b (block localization), together with the inter-block correlation
structure ($r > 0.97$ for all block pairs in both models), indicates that the five
standpoint projections share a dominant global mode. The block-specific holonomy norms
$\|F_{ijk}\|_k$ are not measurements of five independent cognitive processes; they are
projections of a single global coherence signal onto five interpretable axes. This is
consistent with the superposition hypothesis in mechanistic
interpretability~\cite{elhage2022superposition}: transformer representations encode many
features in high-dimensional superposition rather than in cleanly separated subspaces,
making localized block-specific measurements an approximation of a fundamentally
distributed signal. The theoretical framework of Section~\ref{sec:geometric-framework}
remains valid --- the fiber bundle decomposition and gauge invariance results hold
regardless of whether the subspaces are functionally independent. What changes is the
interpretation: we reframe the five standpoint layers from ontologically independent
modules to a diagnostic projection coordinate system
(Section~\ref{sec:five-layer}). Under this reframing, the primary empirical contribution
of the holonomy measure is as a global coherence signal, and the five-axis decomposition
provides structured diagnostic vocabulary for exploring the geometry of this signal.

The empirical evidence warrants a principled reframing of the LCESA framework. The
original theoretical vision --- that transformer internals organize into five
independently measurable standpoint dimensions --- is not confirmed by the data. What the
data do support is a different and arguably more interesting structure: a single dominant
coherence manifold with interpretable projection axes.

Concretely, a principal component analysis of the block holonomy vectors
$(\|F\|_{\min}, \|F\|_{\mathrm{nar}}, \|F\|_{\mathrm{soc}}, \|F\|_{\mathrm{mor}},
\|F\|_{\mathrm{pos}})$ across all 180 Llama-7b test conversations (30 per scenario type;
see Appendix~\ref{app:supplementary}) reveals that the first
principal component accounts for $99.99\%$ of total variance. The PC1 loadings are
approximately uniform across all five standpoint blocks ($\ell_k \approx -0.447 \approx
-1/\sqrt{5}$ for each $k \in K$), confirming that PC1 captures a global coherence signal
rather than any single standpoint dimension. PC1 alone achieves $\varepsilon^2 = 0.42$
for scenario discrimination (Kruskal--Wallallis), matching the full five-dimensional
result ($\varepsilon^2 = 0.42$--$0.68$). This first component is the ``global
complexity mode'' --- a scalar summary of how much the attention transport deviates from
expected paths, regardless of which standpoint layer is examined. The residual components
(PC2--PC5, collectively accounting for $< 0.01\%$ of variance) carry no meaningful
additional diagnostic information beyond the global mode.

We interpret this structure as follows. Transformer attention does not route cognitive
states through five independent channels. Instead, cognitive state differences manifest as
distributed perturbations across the entire residual stream geometry, which our five-axis
projection system captures from five psychologically motivated viewing angles. The moral
block happens to have the highest projection weight in the global mode --- not because
moral processing is uniquely active across all scenarios, but because the value subspaces
assigned to the moral standpoint heads happen to capture a larger fraction of the global
coherence variance in both models tested.

This reframing does not undermine the framework's value; it redirects it. The holonomy
measure $\|F_{ijk}\|$ is a valid and robust global diagnostic. The five-axis projection
provides interpretable vocabulary. Together they constitute a measurement system for
transformer internal geometry, not a decomposition of independent cognitive processes.

% ------------------------------------------------------------------------------
\subsection{Relation to Existing Alignment Methods}

LCESA does not replace existing alignment methods; it provides the internal structure
that makes them more effective.

\paragraph{RLHF.}
Reinforcement learning from human feedback~\cite{ouyang2022training} shapes behavior
through external reward signals. Holonomy minimization is a complementary
\emph{internal} objective: it ensures that commitments made under RLHF training are
preserved across events. The two can be combined: $L_{\mathrm{total}} =
L_{\mathrm{RLHF}} + \lambda \cdot F_{\mathrm{NSI}}$.

\paragraph{Constitutional AI.}
Constitutional AI~\cite{bai2022constitutional} provides a set of principles that the
system should follow. In the LCESA framework, these principles define the alignment
target $\mathcal{T}$: the expected transport $U_{ik}^{\exp}$ is estimated from
conversations that satisfy the constitutional principles. Holonomy deviation then measures
the system's deviation from constitutional consistency.

\paragraph{Mechanistic interpretability.}
Linear probing~\cite{belinkov2021probing} answers ``what information does layer $k$
encode?'' Holonomy deviation answers ``is that information preserved across events?'' The
two diagnostics are complementary: a layer can have high probing accuracy (it encodes the
right commitment) yet high holonomy deviation (it fails to transport that commitment
consistently). Under the global-mode regime (Section~\ref{sec:reframing}), holonomy
deviation primarily measures the overall fidelity of transport rather than layer-specific
failure; the block-specific projections provide interpretable secondary vocabulary.

% ------------------------------------------------------------------------------
\subsection{Architectural Vision: Standpoint as a New Primitive}

The transformer~\cite{vaswani2017attention} introduced \textbf{attention} as a new
architectural primitive --- a mechanism for selecting and propagating information. This
paper introduces \textbf{standpoint} as a second primitive --- a mechanism for
accumulating and constraining commitment. The two primitives play complementary roles:

\begin{itemize}
    \item \textbf{Attention} answers: \emph{What information is relevant?} It selects
    tokens, propagates representations, and computes context-dependent relevance.
    \item \textbf{Standpoint} answers: \emph{What commitments constrain my behavior?} It
    accumulates prior outputs as internal constraints, maintains cross-event coherence,
    and provides a reference frame against which new outputs are evaluated.
\end{itemize}

Current AI systems have attention but lack standpoint. A language model that asserts
``I am committed to X'' in one turn has no architectural basis for that commitment to
exert causal force on its subsequent generations. The result is that behavioral
consistency, however sophisticated, remains a surface phenomenon: it is achieved through
pattern matching in the training data, not through internal structure.

The LCESA framework proposes a concrete path from ``passive response'' to ``active
maintenance.'' The low-holonomy attractor $\mathcal{A}_{\mathrm{self}}$
(Definition~\ref{def:attractor}) is not an external loss function imposed by a trainer;
it is an internal geometric property of the system's own standpoint trajectory. A system
that maintains low holonomy deviation is one that \emph{preserves its own commitments} ---
not because it was told to, but because its architecture makes commitment violation
geometrically costly.

This is not a claim about consciousness. It is a claim about architecture: if we want
AI systems that maintain coherent commitments over time, we need a mechanism for
commitment accumulation. Standpoint is that mechanism.

% ------------------------------------------------------------------------------
\subsection{Future Directions}

\paragraph{Baseline robustness and PCA analysis.}
We have verified that the scenario discrimination (H3) and T0 anomaly
(H7) results are robust to baseline selection via bootstrap ensemble analysis (five
random T1 subsets, 20 bootstrap iterations each): H3~$\varepsilon^2$ ranges from
$0.46$ to $0.51$ across subsets, and T0 ranks first in all subsamples. PCA of the
block holonomy matrix confirms that PC1 accounts for $99.99\%$ of variance with
uniform loadings, establishing the global-mode hypothesis quantitatively. These results
are reported in detail in Appendix~C.

\paragraph{Null grouping controls.}
To verify that the learned head grouping $\gamma$ provides non-trivial diagnostic
decomposition, we plan to compare it against random, shuffled, and layer-uniform null
groupings. If the learned grouping does not significantly outperform null groupings, the
five-layer diagnostic coordinate system has limited empirical support.

\paragraph{Activation patching.}
The most important missing experiment is activation patching: replacing specific layer
activations from T2 scenarios into T1 scenarios to test whether holonomy deviation follows
the activation (causal status) or the prompt text (correlate). This requires GPU resources
and is planned for a subsequent computational run.

\paragraph{Standpoint-aware attention.}
The most direct architectural extension is to design attention mechanisms that explicitly
minimize holonomy deviation. This could take the form of a holonomy penalty in the loss
function, or an architectural modification that enforces block-diagonal transport by
construction.

\paragraph{Integration with RLHF.}
The combination $L_{\mathrm{RLHF}} + \lambda \cdot F_{\mathrm{NSI}}$ is a natural next
step. The key question is whether holonomy regularization improves alignment robustness
--- whether a system trained with both objectives maintains its commitments more
reliably under adversarial pressure.

\paragraph{Multi-agent standpoint.}
The framework extends naturally to multi-agent settings: each agent has its own standpoint
bundle, and inter-agent communication defines transport between bundles. The holonomy
deviation of inter-agent transport measures the coherence of the multi-agent system.
